\section{Introduction}
\label{sec:intro}

In October 2026 I ran a pre-registered gate for a quantization study on a small language model:
a candidate format passed only if its MMLU score stayed at or above 51.0 on 200 questions. It
scored 50.5 and failed by one question. The same report noted that running another format twice,
with identical code, had given 107 and 106 correct MMLU answers, because the evaluation used a
four-slot \texttt{llama-server} whose batches form differently on every run. The gate stayed
failed as written, but it left a question I could not answer at the time: how much of a
measured difference between two quantized models, run locally on a CPU, comes from the models
and how much from how the arithmetic happened to be executed?

This paper measures that. I hold the release, the weights file, the prompts and greedy decoding
fixed and change only settings that the documentation presents as performance options with no
effect on the model: the CPU backend tier the runtime loads, the thread count, the micro-batch
size, flash attention, weight repacking, the server's prompt cache and its number of parallel
slots. I call these \emph{execution-equivalent} configurations. In exact arithmetic they compute
the same function. In float32 on a CPU, most of them do not, and the differences are large
enough to change answers.

The setting is llama.cpp on x86 CPUs, the common way to run quantized GGUF models on consumer
hardware. Two properties make it a good case. First, one prebuilt release ships several
compiled copies of its CPU kernels and picks one at start-up from the machine's instruction-set
flags, so two people running the same command on different laptops silently run different
code. Second, quantized matrix products round activations to 8-bit integers before every
multiplication. Prior work on inference nondeterminism concentrates on GPU serving stacks
\citep{he2025,yuan2025,gond2026,zhao2026,cooper2026}; this paper looks at the CPU path, where
neither batch invariance nor cross-machine agreement has been measured as far as I can find.

\paragraph{Findings.}
On Qwen3-0.6B in Q4\_K\_M, with 980 items from MMLU, TriviaQA and GSM8K:
\begin{itemize}
\item Running the same configuration twice is bitwise reproducible, and thread count never
mattered in any test. Neither did the \texttt{haswell} tier or micro-batches of 8 and 64. Every
other setting changed the logits for at least one file, and the AVX and SSE tiers, micro-batch 1
and flash attention changed them for all six.
\item Running the same file on the AVX-only backend tier instead of the AVX2 tier left no item
bitwise identical, changed 20\% of generated outputs and 41 of 600 MMLU answers, and produced 37
correct/incorrect flips, while every task's accuracy moved by at most 1.3 points. On wikitext the
tier change alone gives a mean KL divergence of 0.011 from the reference and changes the top-1
token at 6\% of positions.
\item The common parallel evaluation setup (four slots, four concurrent clients, prompt cache on)
is not reproducible from run to run: two runs of the identical command agreed bitwise on only
45\% of items and produced different GSM8K reasoning chains for 47 of 80 problems.
\item The size of these effects is explained by activation quantization. With F32 weights and
an F32 KV cache, the same tier change leaves the logits within the measurement floor and changes
no top-1 token; with Q8\_0 weights the mean KL divergence is above $10^{-3}$, and with Q4\_K\_M
above $10^{-2}$. Rounding activations to integers turns ulp-level float differences into
full-step integer differences that propagate through the network.
\end{itemize}

\paragraph{Contributions.}
(1) A measurement of how much execution-equivalent settings change outputs and benchmark scores
for quantized models on CPUs, across six model files and up to twelve settings, with all raw
per-item outputs released. (2) A pre-registered test that identifies activation quantization as
the amplifier, plus code-level explanations, with pre-registered predictions, of which settings
are safe and which are not. (3) A short protocol for local evaluations: which settings to pin,
what to report, and how to size the noise floor before comparing two quantized models.
Section~\ref{sec:background} describes the runtime, Section~\ref{sec:method} the design,
Section~\ref{sec:results} the measurements, Section~\ref{sec:mechanism} the mechanism and
Section~\ref{sec:recommend} the protocol.
